\documentstyle[%


righttag%


,11pt%


,amssymb%


,russian%               remove in Eng.


%,izvru%                 Set izven% in Eng.


,izvru-ks%             for brief communications


]{amsart}





% 





\newcommand{\im}{\operatorname{Im}}


\newcommand{\re}{\operatorname{Re}}


\newcommand{\tts}[1]{}





\theoremstyle{definition}


\theorembodyfont{\rm}


\newtheorem{cond}{\hskip\parindent Условие}





%\hoffset=-15mm%                  For Eng.


\setcounter{page}{73}


\god{2004}


\nomer{3~(502)} %                        Remove in Eng.


%\nomer{Vol.\,48, No.\,3}%               Set in Eng.


%\str{\pageref{begin}--\pageref{end}}%  Set in Eng.


\udk{517.983}





\author{\it М.Ю.\,Кокурин, В.В.\,Ключев}


% NB:   ^^ remove in Eng.


\title{О логарифмических оценках скорости сходимости 


методов решения обратной задачи Коши в банаховом пространстве}








\date{}


%\pagestyle{myheadings}%         Set in Eng.


\pagestyle{plain} %              Remove in Eng.


\begin{document}


%\thispagestyle{plain}%          Set in Eng.


\maketitle


\label{begin}











1. Оценкам скорости сходимости методов решения линейных


некорректных операторных уравнений


\begin{equation}


Bu=f, \quad u\in X,


\end{equation}


посвящено значительное число публикаций (см., напр., [1]--[3] и


цитированную там литературу).


Предположим, что $B$ --- линейный непрерывный оператор,


действующий в


банаховом пространстве $X$ и не обладающий непрерывным обратным.


Ввиду некорректности задачи (1) получение квалифицированных


оценок скорости сходимости методов аппроксимации ее решения возможно лишь


при наложении на искомое решение тех


или иных дополнительных условий. Как правило, эти условия имеют вид


требования истокообразной представимости начальной невязки


$u^* - \xi$, где $u^*$ --- решение (1), $\xi$ --- начальное


приближение, используемое в рассматриваемом методе.





Значительная часть известных методов аппроксимации решений


уравнений (1) может быть получена в рамках общей схемы ([3], с.\,39)


\begin{equation}


u_{\alpha} = (E - \Theta(B,\alpha) B)\xi + \Theta(B,\alpha)f,


\quad \alpha\in (0, \alpha_0],


\end{equation}


определяющей такие приближения $u_\alpha$, что $\lim\limits_{\alpha\to 0}


u_{\alpha} = u^*$. Функция оператора $B$ в (2) понимается в смысле


подходящего операторного исчисления.


В случае секториального оператора $B$ для широкого класса порождающих функций


$\Theta(\lambda, \alpha)$ требование степенной истокопредставимости


$u^* - \xi \in R(B^p)$, $p>0$, 


обеспечивает выполнение степенной оценки скорости сходимости


$\|u_{\alpha}-u^*\|\le C_1{\alpha}^p$ \; $\forall\alpha\in(0, \alpha_0]$


с тем же показателем $p$ (см. [3], с.\,42). Здесь и далее $C_1, C_2,\dotsc$


--- положительные абсолютные константы. В типичных случаях условие 


истокопредставимости означает повышенную гладкость невязки $u^* - \xi$ по


сравнению с гладкостью элементов исходного пространства $X$. В случае 


пространства суммируемых функций $X=L_r$, $1<r<\infty$, это условие


зачастую сводится к требованию


принадлежности начальной невязки $u^* - \xi$ некоторому


пространству Соболева ([3], с.\,24). Если $X$ --- гильбертово


пространство, $B^*=B\geq0$, то аналогичную интерпретацию допускает 


логарифмическое


условие истокопредставимости $u^* - \xi \in R((-\ln B)^{-p})$, $p>0$,


которое влечет логарифмическую оценку скорости сходимости [4]


\begin{equation}


\| u_{\alpha} - u^* \| \leq C_2 (- \ln \alpha)^{-p}


\quad\forall\alpha\in(0,\alpha_0].


\end{equation}


Представляет интерес анализ условий


истокопредставимости, обеспечивающих выполнение логарифмической


оценки скорости сходимости (3) и в случае банахова пространства $X$.





В данной работе рассматривается уравнение (1) с оператором


$B=U(T)$, где $U(t)$, $0\leq t\leq T$, ---


полугруппа ограниченных операторов, действующих в банаховом


пространстве $X$. Устанавливается, что условие истокопредставимости


$u^* - \xi \in R(A^{-p})$, $p>0$,


в котором $- A$ есть генератор полугруппы $U(t)$,


обеспечивает выполнение оценки (3).


Последнее условие может рассматриваться в качестве банахова


аналога требования логарифмической истокопредставимости в


применении к уравнению (1) с оператором $B$ указанного выше вида.


\medskip





2. Уточним постановку задачи. Рассмотрим задачу Коши для абстрактного 


параболического уравнения


\begin{equation}


\frac{dx}{dt} + Ax=0,\quad x(0) = x_0.


\end{equation}


Здесь $A: D(A)\subset X\to X$ --- замкнутый в общем случае неограниченный


оператор с плотной в $X$ областью определения $\overline{D(A)} =


X$. Под решением задачи (4) понимается функция $x=x(t)\in X$, $t\in[0, T]$,


дифференцируемая и удовлетворяющая дифференциальному уравнению (4) при $t>0$ и


непрерывная в точке $t=0$ ([5], c.\,76). Далее через


$\sigma(A)$ обозначается спектр оператора $A$, 


$R(\lambda, A)=(\lambda E - A)^{-1}$ есть резольвента $A$.


Обозначим также $K(\varphi)=\{\lambda\in{\bold C}\setminus\{0\}: 


|\arg\lambda| \leq \varphi \}\cup\{0\}$, $\varphi\in [0, \pi]$; 


$S(R)=\{\lambda\in{\bold C}: |\lambda| \leq R \}$, 


$K(\varphi, R) = K(\varphi)\cap S(R)$.


Известно, что при выполнении условий


\begin{equation}


\sigma(A)\subset K(\varphi_0),\quad


\| R(\lambda, A)\| \leq C_3/|\lambda|


\quad\forall\lambda\in{\bold C}\setminus K(\varphi_0),


\quad \varphi_0\in (0,\pi/2),


\end{equation}


задача Коши (4) является корректной ([5], с.\,38).


В этом случае решение (4) имеет вид $x = U(t) x_0$, где $U(t)$


--- сильно непрерывная полугруппа линейных ограниченных операторов 


в $X$, генератор которой совпадает с оператором $-A$.


В предположениях (5) обратная задача Коши


для уравнения $dx/dt +Ax = 0$, состоящая в определении


начального условия $x_0=x(0)$ по заданному состоянию $x(T)$ в момент $T>0$,


является в общем случае некорректной ([5], c.\,101). Положив


$B=U(T)$, $x_0=u$, $x(T)=f$, представим эту задачу в виде (1).





В дальнейшем считаем выполненным основное 





\begin{cond}


Оператор $A$ обладает компактным обратным $A^{-1}:X\to X$, спектр


$\sigma(A)$ принадлежит полуполосе


$$


\Pi = \{\lambda\in{\bold C}: | \im


\lambda | \leq \psi_0/T,\ \re\lambda \geq a \}\subset K(\varphi_0),


$$


где $\psi_0\in(0, \pi/2)$, $a>0$, и выполняется оценка из (5).


\end{cond}





Относительно порождающей функции $\Theta(z, \alpha)$ предположим, что она 


удовлетворяет следующим условиям.





\begin{cond}


Функция $\Theta(z,\alpha)$ аналитична


по $z$ в открытой окрестности $D_{\alpha}$ множества


$K(\psi_0, 1)\cup S(\alpha/2)$, $\alpha\in(0,\alpha_0]$.


\end{cond}





\begin{cond}


При $\alpha\in(0,\alpha_0]$ имеет место соотношение


$$


\int_{\Gamma_0(\alpha)} | 1 - \Theta(z,\alpha)z |


| dz | \leq C_4\alpha,


$$


где контур $\Gamma_0(\alpha)$ окружает множество


$K(\psi_0, 1)\cup S(\alpha/2)$. 


\end{cond}





Условиям 2, 3 удовлетворяют многие распространенные в вычислительной практике


порождающие функции $\Theta(z, \alpha)$. В качестве примера рассмотрим функцию


\begin{equation}


\Theta (z,\alpha) =


z^{-1}\big(1 - (1-\mu_0 z)^{1/\alpha}\big),


\quad\alpha=1/n,\quad n\in {\bold N}.


\end{equation}


Схема (2), (6) сводится к простейшему итерационному процессу ([3], c.\,77) 


$u_\alpha=u^{(n)}$, где


$$


u^{(0)} = \xi, \quad


u^{(k+1)} = u^{(k)} - \mu_0(\exp(-TA)u^{(k)} - f),


\quad k =0, 1,\dotsc, n-1.


$$


Здесь для нахождения элемента $v^{(k)}=\exp(-TA)u^{(k)}=U(T)u^{(k)}$ решается


корректная задача Коши


$$


\frac{dx}{dt} + Ax = 0,\quad x(0) = u^{(k)}.


$$


Функция (6) удовлетворяет условиям 2 и 3,


если в качестве $\Gamma_0(\alpha)$ выбрать границу множества 


$K(\psi_0, 1)\cup S(\alpha/2)$ из условия 2.





Другие примеры порождающих функций, удовлетворяющие условиям 2, 3,


можно найти, например, в ([3], c.\,64--76).





\begin{thm}


Пусть выполняются условия $1$--$3$


и истокообразное представление $u^* - \xi \in R(A^{-p})$, ${p>0}$.


Тогда для приближений $u_\alpha$,


определяемых согласно $(2)$ при $B=\exp(-TA)$, имеет место оценка $(3)$.


\end{thm}





По схеме (2) эффективно строятся алгоритмы регуляризации 


уравнения (1) в случае, когда правая часть $f=x(T)$


известна с погрешностью, т.\,е. задан элемент $f_\delta$ такой,


что $\| f_\delta - f \| \leq \delta$, $\delta>0$.


В соответствии с ([3], с.\,43) в качестве приближения к решению уравнения


(1) в этом случае выбирается элемент


\begin{equation}


u_\alpha^\delta = (E - \Theta(B,\alpha)B)\xi +


\Theta(B,\alpha) f_\delta.


\end{equation}


Параметр регуляризации $\alpha=\alpha(\delta)$ в (7)


следует согласовать с погрешностью $\delta$ так, чтобы


выполнялись соотношения


$\lim\limits_{\delta\to 0} \alpha(\delta) = 0$ и


$\lim\limits_{\delta\to 0}


\sup\{ \| u_{\alpha(\delta)}^\delta-u^*\|:\|f_\delta-f\|\leq\delta\}=0$


([2], с.\,27).





Наложим на порождающую функцию $\Theta(z,\alpha)$ дополнительное





\begin{cond}


Для любых $\alpha\in(0,\alpha_0]$


имеет место оценка


$$


\int_{\Gamma_0(\alpha)} | z|^{-1}|\Theta(z,\alpha)|\,


| d z| \leq C_5{\alpha}^{-1}.


$$


\end{cond}





\begin{thm}


Пусть выполняются условия $1$--$4$ и истокообразное представление 


$u^* - \xi \in R(A^{-p})$, ${p>0}$.


Тогда при $\alpha(\delta)=C_6{\delta}^s$, $s\in(0, 1)$, 


для приближений, вырабатываемых процессом $(7)$,


при достаточно малых $\delta>0$ имеет место оценка


$$


\| u_{\alpha(\delta)}^{\delta} - u^* \| \leq C_7(-\ln\delta)^{-p}.


$$


\end{thm}





Другие подходы к построению методов регуляризации рассматриваемой 


обратной задачи можно найти, например, в [6]. Можно показать, что в случае, 


когда оператор $A$ определяется регулярно эллиптическим


дифференциальным выражением, условие $u^* - \xi \in R(A^{-p})$, $p>0$, 


эквивалентно требованию принадлежности начальной невязки 


соответствующему пространству Соболева--Бесова. 








\begin{thebibliography}{9}





\bibitem{1}


Иванов В.К., Васин В.В., Танана В.П. {\em Теория линейных некорректных 


задач и ее приложения}. -- М.: Наука, 1978. -- 208~с.


\bibitem{2}


Бакушинский А.Б., Гончарский А.В. {\em Итеративные методы решения 


некорректных задач}. -- М.: Наука, 1989.-- 128~с. 


\bibitem{3}


Бакушинский А.Б., Кокурин М.Ю. {\em Итерационные методы решения некорректных 


операторных уравнений с гладкими операторами}. -- М.: Едиториал УРСС, 


2002. -- 192~с.


\bibitem{4}


Hohage T. {\em Regularization of exponentially ill-posed problems}


// Numer. Funct. Anal. and Optimization. -- 2000. -- V.\,21.


-- \No\No\,3,~4. -- P.\,439--464.


\bibitem{5}


Крейн С.Г. {\em Линейные дифференциальные уравнения в банаховом пространстве}.


-- М.: Наука, 1967. -- 464~с.


\bibitem{6}


Иванов В.К., Мельникова И.В., Филинков А.И. {\em Дифференциально-операторные


уравнения и некорректные задачи}. -- М.: Физматлит, 1995. -- 176~с.





\end{thebibliography}





\vskip 1cm





\post{\it Марийский государственный}{\it Поступила}


\post{\it университет}{08.01.2003}





\label{end}


\end{document}


